\documentclass[12pt]{article}
\usepackage[utf8]{inputenc}
\usepackage[T1]{fontenc}
\usepackage[margin=1in]{geometry}
\usepackage{setspace}
\usepackage{xcolor}
\usepackage[colorlinks=true, urlcolor=blue!60!black, linkcolor=blue!60!black]{hyperref}
\setlength{\parindent}{2em}
\setlength{\parskip}{1.25em}
\renewcommand{\baselinestretch}{1.0}
\begin{document}
\begin{center}
{\LARGE \textit{United States v. Skrmetti}, 605 U.S. 495 (2025)\par}
\vspace{0.8em}
{\large Case Brief\par}
\end{center}
\vspace{1em}

\section{Memory Jogger}

Tennessee's ban on puberty blockers and hormone therapy for minors to treat gender dysphoria classifies by age and medical use rather than sex or transgender status, so it triggers only rational basis review under the Equal Protection Clause and is constitutional (\href{https://www.courtlistener.com/opinion/10776831/united-states-v-skrmetti/}{CourtListener}; \href{https://www.law.cornell.edu/supremecourt/text/23-477}{Cornell LII}).

\section{Detailed Case Facts}

In March 2023, Tennessee joined more than 20 states restricting sex-transition medical interventions for minors by enacting the Prohibition on Medical Procedures Performed on Minors Related to Sexual Identity, Senate Bill 1 (SB1), codified at Tenn. Code Ann. \S 68-33-101 \textit{et seq.} (\href{https://www.law.cornell.edu/supremecourt/text/23-477}{Cornell LII}; \href{https://www.supremecourt.gov/opinions/24pdf/23-477_2cp3.pdf}{Slip Opinion}). SB1 prohibits healthcare providers from surgically modifying tissues, cavities, or organs or from prescribing, administering, or dispensing puberty blockers or hormones to any minor for either of two purposes: (1) enabling a minor to identify with, or live as, a purported identity inconsistent with the minor's biological sex, or (2) treating purported discomfort or distress from a discordance between the minor's biological sex and asserted identity (Tenn. Code Ann. \S\S 68-33-102(5), 68-33-103(a)(1)) (\href{https://www.law.cornell.edu/supremecourt/text/23-477}{Cornell LII}).

SB1 contains two key statutory limits. First, it applies only to minors under age 18 and does not restrict adults (\S 68-33-102(6)). Second, it permits healthcare providers to prescribe puberty blockers and hormones to minors to treat a congenital defect, precocious puberty, disease, or physical injury (\S 68-33-103(b)(1)(A)), while explicitly excluding ``gender dysphoria, gender identity disorder, [and] gender incongruence'' from the definitions of ``congenital defect'' and ``disease'' (\S\S 68-33-102(1), 68-33-103(b)(2)) (\href{https://www.law.cornell.edu/supremecourt/text/23-477}{Cornell LII}). The statute defines ``sex'' as a person's immutable reproductive characteristics determined by anatomy and genetics at birth (\S 68-33-102(9)) (\href{https://harvardlawreview.org/print/vol-139/skrmetti-beyond-scrutiny/}{Harvard Law Review}).

In its legislative findings, the Tennessee General Assembly stated that the prohibited interventions can cause irreversible sterility, increased risk of disease, and adverse psychological outcomes; that minors lack the maturity to appreciate these life-altering consequences and may later experience regret; that European health authorities in Sweden, Finland, and the United Kingdom had restricted such treatments due to weak evidence; and that the State has a compelling interest in ``encouraging minors to appreciate their sex'' and prohibiting care ``that might encourage minors to become disdainful of their sex'' (\S 68-33-101(b)--(e), (h), (m)) (\href{https://www.law.cornell.edu/supremecourt/text/23-477}{Cornell LII}). SB1 is enforced through three mechanisms: Tennessee Attorney General enforcement actions seeking injunctions, profit disgorgement, and civil penalties of \$25,000 per violation (\S 68-33-106(b)); professional licensing discipline (\S 68-33-107); and a private right of action for an injured minor or nonconsenting parent (\S 68-33-105) (\href{https://www.law.cornell.edu/supremecourt/text/23-477}{Cornell LII}).

Three transgender adolescents in Tennessee (L.W., John Doe, and Ryan Roe), their parents, and a Tennessee physician filed a pre-enforcement lawsuit against Tennessee Attorney General Jonathan Skrmetti and state officials, arguing that SB1 violates the Equal Protection Clause of the Fourteenth Amendment (and, in the lower courts, substantive due process parental rights) (\href{https://www.law.cornell.edu/supremecourt/text/23-477}{Cornell LII}; \href{https://www.oyez.org/cases/2024/23-477}{Oyez}). The United States intervened under 42 U.S.C. \S 2000h-2 after the Attorney General certified the case as one of general public importance (\href{https://www.law.cornell.edu/supremecourt/text/23-477}{Cornell LII}). Following the change in presidential administrations after oral argument, the Department of Justice sent a letter to the Supreme Court on February 7, 2025, reversing its position and stating that SB1 does not violate the Equal Protection Clause, while declining to seek dismissal because the private plaintiffs remained adverse to Tennessee (\href{https://www.law.cornell.edu/supremecourt/text/23-477}{Cornell LII}).

\section{Procedural History}

\begin{itemize}
\item \textbf{U.S. District Court for the Middle District of Tennessee:} In \textit{L.W. v. Skrmetti}, 679 F. Supp. 3d 668 (M.D. Tenn. 2023), the district court held that the plaintiffs lacked standing to challenge SB1's surgical ban, but were likely to succeed on their Equal Protection challenge to the ban on puberty blockers and hormones. Concluding that transgender persons constitute a quasi-suspect class, that SB1 discriminates based on sex and transgender status, and that SB1 likely failed intermediate scrutiny, the court issued a statewide preliminary injunction barring enforcement of SB1's medication prohibitions (while leaving the surgical ban and private right of action undisturbed) (\href{https://www.law.cornell.edu/supremecourt/text/23-477}{Cornell LII}).
\item \textbf{U.S. Court of Appeals for the Sixth Circuit:} A divided panel first stayed the preliminary injunction pending appeal (\textit{L.W. v. Skrmetti}, 73 F.4th 408 (6th Cir. 2023)) and then reversed in \textit{L.W. v. Skrmetti}, 83 F.4th 460 (6th Cir. 2023) (opinion by Chief Judge Sutton; dissent by Judge White) (\href{https://www.law.cornell.edu/supremecourt/text/23-477}{Cornell LII}; \href{https://ballotpedia.org/United_States_v._Skrmetti}{Ballotpedia}). The Sixth Circuit held that SB1 regulates treatments for all minors regardless of sex, declined to recognize transgender individuals as a suspect class, found no animus, and upheld SB1 under rational basis review. Judge White dissented, arguing that SB1 facially discriminates based on sex assigned at birth and fails intermediate scrutiny (\href{https://www.law.cornell.edu/supremecourt/text/23-477}{Cornell LII}).
\item \textbf{U.S. Supreme Court:} The United States petitioned for certiorari on November 6, 2023. The Court granted certiorari on June 24, 2024, limited to the Equal Protection Clause question (declining to grant the private plaintiffs' separate petition raising a substantive due process parental-rights claim), heard oral argument on December 4, 2024, and affirmed the Sixth Circuit on June 18, 2025 (\href{https://www.scotusblog.com/cases/united-states-v-skrmetti/}{SCOTUSblog}; \href{https://ballotpedia.org/United_States_v._Skrmetti}{Ballotpedia}).
\end{itemize}

\section{Judicial Votes}

\textbf{6--3 decision} affirming the Sixth Circuit (\href{https://www.law.cornell.edu/supremecourt/text/23-477}{Cornell LII}; \href{https://www.scotusblog.com/cases/united-states-v-skrmetti/}{SCOTUSblog}):

\begin{itemize}
\item \textbf{Majority Opinion (Chief Justice Roberts):} Joined in full by Justices Thomas, Gorsuch, Kavanaugh, and Barrett; joined by Justice Alito as to Parts I and II-B.
\item \textbf{Concurring Opinion (Justice Thomas):} Joined the Court's opinion in full and wrote separately.
\item \textbf{Concurring Opinion (Justice Barrett):} Joined the Court's opinion in full and wrote separately; joined by Justice Thomas.
\item \textbf{Concurring in Part and Concurring in the Judgment (Justice Alito):} Joined Parts I and II-B of the Court's opinion, declined to join Part II-A (Parts II-A-1, II-A-2, and II-A-3), and concurred in the judgment.
\item \textbf{Dissenting Opinion (Justice Sotomayor):} Joined in full by Justice Jackson; joined by Justice Kagan as to Parts I--IV (all except Part V).
\item \textbf{Dissenting Opinion (Justice Kagan):} Wrote separately to explain why she joined Parts I--IV of Justice Sotomayor's dissent but not Part V.
\end{itemize}

\section{Holding}

\textbf{Judgment:} Affirmed (6--3). The judgment of the U.S. Court of Appeals for the Sixth Circuit reversing the district court's preliminary injunction is affirmed (\href{https://www.law.cornell.edu/supremecourt/text/23-477}{Cornell LII}).

\textbf{Rule Announced:} A state law prohibiting healthcare providers from prescribing or administering puberty blockers or hormones to minors to treat gender dysphoria, gender identity disorder, or gender incongruence---while permitting those medications for other medical diagnoses---classifies on the basis of age and medical use rather than sex or transgender status. Because neither classification triggers heightened scrutiny under the Equal Protection Clause of the Fourteenth Amendment, such a law is reviewed under rational basis review and is constitutional so long as it is rationally related to a legitimate state interest, such as protecting minors' health and welfare amid medical and scientific uncertainty (\href{https://www.law.cornell.edu/supremecourt/text/23-477}{Cornell LII}; \href{https://supreme.justia.com/cases/federal/us/605/23-477/}{Justia}).

\textbf{Holding vs. Dicta:}
\begin{itemize}
\item \textbf{Holding:} (1) SB1's prohibitions on puberty blockers and hormones for minors do not facially classify on the basis of sex or transgender status; (2) even assuming \textit{Bostock v. Clayton County}, 590 U.S. 644 (2020), applied beyond Title VII, neither sex nor transgender status is a but-for cause of SB1's operation; and (3) SB1 satisfies rational basis review under the Equal Protection Clause (\href{https://www.law.cornell.edu/supremecourt/text/23-477}{Cornell LII}).
\item \textbf{Expressly Left Open / Dicta:} The Court \textit{did not decide} whether transgender individuals constitute a suspect or quasi-suspect class (only Justices Barrett, Thomas, and Alito addressed and rejected suspect-class status in separate concurrences); \textit{did not decide} whether \textit{Bostock}'s Title VII reasoning extends to the Equal Protection Clause; \textit{did not address} SB1's ban on sex-transition surgery (for which plaintiffs lacked standing) or any substantive due process claim regarding parental rights; and noted in footnote 3 that its analysis does not govern a law regulating a class of persons defined by a sex-linked characteristic rather than a class of medical treatments or conditions (\href{https://www.law.cornell.edu/supremecourt/text/23-477}{Cornell LII}; \href{https://fedsoc.org/commentary/fedsoc-blog/supreme-court-upholds-tennessee-s-ban-on-gender-affirming-care-for-minors-in-united-states-v-skrmetti}{Federalist Society}).
\end{itemize}

\section{Analysis of Opinions}

\subsection{Majority Opinion (Chief Justice Roberts)}

Chief Justice Roberts held that SB1 does not trigger heightened scrutiny because on its face it draws only two classifications---age (minors versus adults) and medical use (permissible diagnoses such as precocious puberty or congenital defects versus excluded diagnoses of gender dysphoria, gender identity disorder, and gender incongruence)---both of which receive rational basis review under \textit{Massachusetts Board of Retirement v. Murgia}, 427 U.S. 307 (1976), and \textit{Vacco v. Quill}, 521 U.S. 793 (1997) (\href{https://www.law.cornell.edu/supremecourt/text/23-477}{Cornell LII}).

First, Roberts rejected the claim that SB1 classifies by sex (Part II-A-1). Mere statutory reference to sex does not trigger intermediate scrutiny, especially in the medical context where treatments and FDA approvals are often bound up with biological sex. Properly understood, a ``medical treatment'' consists of both the drug and the medical indication it treats; thus, a transgender boy taking puberty blockers for gender dysphoria receives a different medical treatment than a boy taking puberty blockers for precocious puberty. Unlike the antimiscegenation statute in \textit{Loving v. Virginia}, 388 U.S. 1 (1967), SB1 does not prohibit conduct for one sex that it allows for the other. Nor does SB1's reference to ``encouraging minors to appreciate their sex'' establish impermissible sex stereotyping, absent a showing of invidious discriminatory purpose under \textit{Personnel Administrator of Massachusetts v. Feeney}, 442 U.S. 256 (1979) (\href{https://www.law.cornell.edu/supremecourt/text/23-477}{Cornell LII}).

Second, relying on \textit{Geduldig v. Aiello}, 417 U.S. 484 (1974), Roberts held that SB1 does not classify based on transgender status (Part II-A-2). Just as California's exclusion of pregnancy from disability coverage divided recipients into ``pregnant women and nonpregnant persons,'' SB1 divides minors into those seeking puberty blockers or hormones for excluded diagnoses (only transgender minors) and those seeking them for other conditions (both transgender and nontransgender minors). Because of this ``lack of identity'' and the absence of any pretext argument, the Court declined to decide whether transgender status is a suspect or quasi-suspect class (\href{https://www.law.cornell.edu/supremecourt/text/23-477}{Cornell LII}).

Third, Roberts concluded that \textit{Bostock v. Clayton County} did not alter the outcome (Part II-A-3). Without deciding whether \textit{Bostock} applies outside Title VII, Roberts reasoned that changing a transgender boy's biological sex from female to male would not allow him to obtain testosterone under SB1 because he would still lack a qualifying diagnosis such as a congenital defect or disease. Unlike sexual orientation in \textit{Bostock}, changing a patient's sex does not automatically change the medical diagnosis (\href{https://www.law.cornell.edu/supremecourt/text/23-477}{Cornell LII}).

Finally, applying rational basis review (Part II-B), Roberts held that SB1 easily survives. Tennessee cited risks of irreversible sterility, health complications, adolescent immaturity, regret, and evolving international medical evidence---including the 2024 Cass Review in England and restrictions in Sweden, Finland, and Norway. Under \textit{Gonzales v. Carhart}, 550 U.S. 124 (2007), states possess wide discretion to legislate amid medical and scientific uncertainty, leaving policy debates to the democratic process (\href{https://www.law.cornell.edu/supremecourt/text/23-477}{Cornell LII}).

\subsection{Concurring Opinion (Justice Thomas)}

Justice Thomas joined the Court's opinion in full and wrote separately to make two broader points (\href{https://www.law.cornell.edu/supremecourt/text/23-477}{Cornell LII}). First, he argued that \textit{Bostock}'s Title VII textual analysis should never be imported into the Equal Protection Clause, which lacks Title VII's ``because of'' statutory wording; extending \textit{Bostock} to constitutional law would unsettle sex-specific medical regulations and threaten equal protection doctrine across other protected traits. Second, Thomas forcefully challenged deference to claimed ``medical consensus.'' Reviewing the off-label status and risks of puberty blockers, cross-sex hormones, and surgeries, along with European systematic reviews, detransitioner testimony, and revelations that WPATH adjusted its 2022 Standards of Care in response to legal strategy and political pressure, Thomas warned that courts must defer to democratic legislatures rather than ``self-described experts'' (\href{https://www.law.cornell.edu/supremecourt/text/23-477}{Cornell LII}).

\subsection{Concurring Opinion (Justice Barrett, Joined by Justice Thomas)}

Justice Barrett joined the Court's opinion in full and wrote separately to answer the question the majority left open: whether transgender status constitutes a suspect or quasi-suspect class (\href{https://www.law.cornell.edu/supremecourt/text/23-477}{Cornell LII}). Applying the \textit{Carolene Products} and \textit{Lyng v. Castillo} factors, Barrett concluded that transgender individuals do not constitute a suspect class because (1) transgender status is not defined by an ``obvious, immutable, or distinguishing characteristic'' ascertainable at birth (citing varying onset ages and detransition); (2) the population is a ``large, diverse, and amorphous'' umbrella category rather than a discrete group; and (3) recognizing suspect-class status would force courts to micromanage legitimate legislative policy choices in healthcare, restrooms, and athletics. Barrett further proposed that future suspect-class inquiries should require a demonstrated history of \textit{de jure} discrimination---discrimination enshrined in law by state actors---rather than private animus, as a more objective and judicially manageable standard (\href{https://www.law.cornell.edu/supremecourt/text/23-477}{Cornell LII}).

\subsection{Opinion Concurring in Part and Concurring in the Judgment (Justice Alito)}

Justice Alito joined Parts I and II-B of the Court's opinion and concurred in the judgment, but declined to join Part II-A (\href{https://www.law.cornell.edu/supremecourt/text/23-477}{Cornell LII}). On sex classification, Alito agreed that SB1 does not classify by sex because equal protection precedents have always defined ``sex'' as biological male or female status and required an overt rule treating men differently from women, rejecting \textit{Bostock}'s counterfactual test in the constitutional context. However, Alito expressed uneasiness with the majority's \textit{Geduldig} reasoning in Part II-A-2, acknowledging a ``strong argument'' that SB1 does classify on the basis of transgender status because it targets a minor's intent to live inconsistently with biological sex. Rather than avoid the issue, Alito assumed \textit{arguendo} that SB1 classifies by transgender status and concluded---similarly to Justice Barrett---that transgender individuals do not qualify as a suspect or quasi-suspect class, leaving SB1 subject to and valid under rational basis review (\href{https://www.law.cornell.edu/supremecourt/text/23-477}{Cornell LII}).

\subsection{Dissenting Opinion (Justice Sotomayor, Joined by Justice Jackson and in Part by Justice Kagan)}

Justice Sotomayor, joined in full by Justice Jackson and in Parts I--IV by Justice Kagan, argued that SB1 plainly classifies on the basis of sex on its face and therefore requires intermediate scrutiny (\href{https://www.law.cornell.edu/supremecourt/text/23-477}{Cornell LII}). Under SB1, whether a Tennessee doctor may prescribe puberty blockers or hormones depends directly on the adolescent's sex assigned at birth: a biological female with unwanted facial hair (hirsutism) may receive testosterone suppressors, whereas a biological male experiencing identical distress may not; a biological male with benign breast tissue (gynecomastia) may receive hormones to look more masculine, whereas a biological female may not. Analogizing to a law barring minors from attending religious services ``inconsistent with'' their religion, Sotomayor contended that any statute turning on inconsistency with a protected trait classifies on that trait (\href{https://www.law.cornell.edu/supremecourt/text/23-477}{Cornell LII}).

Sotomayor criticized the majority's reliance on \textit{Vacco v. Quill} and \textit{Geduldig}, arguing that SB1's ``medical purpose'' is itself defined by sex and that \textit{Bostock}'s logic confirms sex is at least one but-for cause of SB1's operation. She also argued that transgender individuals satisfy the criteria for at least quasi-suspect class status, pointing to historical and ongoing \textit{de jure} and private discrimination. In Part V (not joined by Justice Kagan), Sotomayor concluded that SB1 would fail intermediate scrutiny given its categorical ban on care endorsed by major American medical associations (\href{https://www.law.cornell.edu/supremecourt/text/23-477}{Cornell LII}; \href{https://ballotpedia.org/United_States_v._Skrmetti}{Ballotpedia}).

\subsection{Dissenting Opinion (Justice Kagan)}

Justice Kagan wrote a brief separate dissent joining Parts I--IV of Justice Sotomayor's opinion to emphasize that SB1 triggers heightened scrutiny as a sex-based classification, which is designed to ``smoke out'' whether a law rests on stereotypes or legitimate state health interests (\href{https://ballotpedia.org/United_States_v._Skrmetti}{Ballotpedia}). However, she declined to join Part V and took ``no view on how SB1 would fare under heightened scrutiny,'' noting that the factual record was extensive and disputed, that the Sixth Circuit had never applied intermediate scrutiny, and that both the United States and the private plaintiffs had asked only for a remand to apply heightened scrutiny in the lower courts (\href{https://ballotpedia.org/United_States_v._Skrmetti}{Ballotpedia}).

\section{Examples: Future Applications}

\begin{enumerate}
\item \textbf{Same-Side (Rational Basis Applies):} A state statute prohibits healthcare providers from prescribing puberty blockers or cross-sex hormones to minors diagnosed with gender dysphoria while permitting those same drugs for endocrine disorders or precocious puberty, and imposes licensing penalties on violators. Under \textit{Skrmetti}, this law classifies by age and medical diagnosis rather than sex or transgender status, triggers only rational basis review, and will be upheld if supported by plausible health-and-safety concerns (\href{https://www.law.cornell.edu/supremecourt/text/23-477}{Cornell LII}).
\item \textbf{Same-Side (Rational Basis Applies):} A state Medicaid program excludes coverage for hormone therapy and surgical procedures prescribed to treat gender dysphoria for beneficiaries of any age, while covering the same hormones and surgeries for cancer, trauma, or congenital abnormalities. Extending \textit{Skrmetti}'s reliance on \textit{Geduldig v. Aiello} and its ``medical use'' classification framework, a court would likely treat this exclusion as a diagnosis-based coverage decision rather than a facial sex or transgender-status classification, subjecting it to rational basis review absent proof of invidious pretext (\href{https://harvardlawreview.org/print/vol-139/skrmetti-beyond-scrutiny/}{Harvard Law Review}).
\item \textbf{Opposite-Side (Heightened Scrutiny Applies):} A state law prohibits public school students from wearing clothing or hairstyles ``inconsistent with'' their biological sex (for example, forbidding biological boys from wearing skirts or long hair while allowing biological girls to do so). Both the majority (Part II-A-1) and Tennessee's own briefing acknowledged that non-medical rules prohibiting sex-inconsistent clothing or occupations directly classify on the basis of sex and enforce sex-based conformity, triggering intermediate scrutiny (\href{https://www.law.cornell.edu/supremecourt/text/23-477}{Cornell LII}).
\item \textbf{Opposite-Side (Heightened Scrutiny Applies):} A state statute denies state-funded health insurance altogether to any individual who ``has ever undergone gender transition'' or ``identifies as transgender,'' barring them from coverage even for unrelated conditions like broken bones or diabetes. Footnote 3 of Chief Justice Roberts's majority opinion explicitly distinguishes a law regulating a class of \textit{persons} defined by a specified characteristic from a law regulating a class of \textit{medical treatments or conditions}, placing such a status-based exclusion outside \textit{Skrmetti}'s \textit{Geduldig} shield (\href{https://www.law.cornell.edu/supremecourt/text/23-477}{Cornell LII}).
\item \textbf{Fence-Sitter (Borderline / Unresolved):} A state law prohibits minors from updating the sex marker on their state-issued birth certificates to reflect a gender identity inconsistent with biological sex, while permitting corrections for clerical errors. On one hand, states may argue under \textit{Skrmetti} and the Barrett/Alito concurrences that birth certificates record historical biological facts and that transgender status is not a suspect class; on the other hand, challengers argue that vital-records rules lie outside \textit{Skrmetti}'s ``medical context'' and directly classify persons by sex and transgender status. Notably, shortly after \textit{Skrmetti}, the Supreme Court granted certiorari, vacated, and remanded a birth-certificate challenge (\textit{Stitt v. Fowler}) for reconsideration in light of \textit{Skrmetti}, highlighting the genuine doctrinal uncertainty (\href{https://harvardlawreview.org/print/vol-139/skrmetti-beyond-scrutiny/}{Harvard Law Review}).
\end{enumerate}

\section{Critique}

\textbf{Originalist and Textualist Perspectives:} Conservative and originalist commentators praised \textit{Skrmetti} for restoring democratic self-governance over contested bioethical questions and halting the expansion of \textit{Bostock}'s Title VII textualism into the Fourteenth Amendment. Writing for the Federalist Society, Sarah Parshall Perry emphasized that the Court properly recognized SB1's age and medical-use classifications, refused to constitutionalize disputed clinical guidelines, and left evolving pediatric-medicine policy to voters and state legislatures (\href{https://fedsoc.org/commentary/fedsoc-blog/supreme-court-upholds-tennessee-s-ban-on-gender-affirming-care-for-minors-in-united-states-v-skrmetti}{Federalist Society}). Scholars studying Justice Barrett's concurrence have also noted that her emphasis on \textit{de jure} discrimination attempts to ground Equal Protection tiers of scrutiny in objective, historically verifiable legal criteria while preserving existing heightened scrutiny for biological sex (\href{https://direct.mit.edu/ajle/article/doi/10.1162/AJLE.a.107/139252/AMY-CONEY-BARRETT-S-SEX-EQUALITY-CONSTITUTION-The}{American Journal of Law and Equality}). Yet even some formalist-leaning observers note internal tension: Justice Alito himself questioned the majority's \textit{Geduldig} move, conceding that SB1 plausibly classifies by transgender status, and some scholars argue that reconciling \textit{Skrmetti} with \textit{Bostock}'s strict textualist but-for test is analytically strained (\href{https://www.law.cornell.edu/supremecourt/text/23-477}{Cornell LII}).

\textbf{Progressive and Doctrinal Perspectives:} Progressive constitutional scholars sharply criticized \textit{Skrmetti} as a doctrinal retreat from anti-stereotyping principles. In the \textit{Harvard Law Review}, Professor Kate Redburn argued that \textit{Skrmetti} creates a troubling ``double deference'' in an undefined ``medical context'': first deferring to the state's framing of a sex-based exclusion as a neutral medical category, and then applying the most deferential rational basis review (\href{https://harvardlawreview.org/print/vol-139/skrmetti-beyond-scrutiny/}{Harvard Law Review}). Critics contend that by reviving \textit{Geduldig v. Aiello} (following \textit{Dobbs v. Jackson Women's Health Organization}), the Court enables legislatures to insulate sex and transgender classifications from heightened scrutiny simply by drafting prohibitions around sex-linked medical diagnoses. Administrative and civil-rights scholars have likewise warned that lower courts risk misapplying \textit{Skrmetti}'s constitutional classification reasoning to narrow statutory Title VII protections under \textit{Bostock} (\href{https://harvardlawreview.org/blog/2026/04/intermediary-choice-and-the-structural-logic-of-title-viis-but-for-test/}{Harvard Law Review Blog}).

\textbf{Underlying Competing Values:} At its core, \textit{Skrmetti} reflects a deep clash between \textit{democratic majoritarianism and legislative flexibility} in areas of scientific uncertainty (especially involving minors' capacity to consent to potentially irreversible treatments) and \textit{countermajoritarian judicial protection} for vulnerable minorities against laws that enforce conformity with birth-assigned sex and override individualized decisions made by parents and physicians.

\section{Key Quotations}

\begin{quote}
``In the medical context, the mere use of sex-based language does not sweep a statute within the reach of heightened scrutiny.''
\end{quote}
\noindent --- Chief Justice Roberts, Opinion of the Court (\href{https://www.law.cornell.edu/supremecourt/text/23-477}{Cornell LII}).

\begin{quote}
``The Equal Protection Clause does not resolve these disagreements. Nor does it afford us license to decide them as we see best. Our role is not `to judge the wisdom, fairness, or logic' of the law before us, \textit{Beach Communications}, 508 U. S., at 313, but only to ensure that it does not violate the equal protection guarantee of the Fourteenth Amendment. Having concluded it does not, we leave questions regarding its policy to the people, their elected representatives, and the democratic process.''
\end{quote}
\noindent --- Chief Justice Roberts, Opinion of the Court (\href{https://www.law.cornell.edu/supremecourt/text/23-477}{Cornell LII}).

\begin{quote}
``This case carries a simple lesson: In politically contentious debates over matters shrouded in scientific uncertainty, courts should not assume that self-described experts are correct.''
\end{quote}
\noindent --- Justice Thomas, concurring (\href{https://www.law.cornell.edu/supremecourt/text/23-477}{Cornell LII}).

\begin{quote}
``The Equal Protection Clause does not demand heightened judicial scrutiny of laws that classify based on transgender status. Rational-basis review applies, which means that courts must give legislatures flexibility to make policy in this area.''
\end{quote}
\noindent --- Justice Barrett, concurring (\href{https://www.law.cornell.edu/supremecourt/text/23-477}{Cornell LII}).

\begin{quote}
``The Court's willingness to do so here does irrevocable damage to the Equal Protection Clause and invites legislatures to engage in discrimination by hiding blatant sex classifications in plain sight. It also authorizes, without second thought, untold harm to transgender children and the parents and families who love them.''
\end{quote}
\noindent --- Justice Sotomayor, dissenting (\href{https://ballotpedia.org/United_States_v._Skrmetti}{Ballotpedia}; \href{https://www.courtlistener.com/opinion/10776831/united-states-v-skrmetti/}{CourtListener}).

\end{document}
